%%%%%%%%%%%%%%%%%%%%%%%%%%%%%%%%%%%%%%%%%%%%%%%%%%%%%%%%%%%%%%%%%%%%%%
% Câu chưa có mã dạng — cần gắn %[..] rồi chuyển vào đúng file dạng
% Số câu: 12  (sinh tự động bởi tach_cau.py; từ nay thầy sửa trực tiếp file này)
%%%%%%%%%%%%%%%%%%%%%%%%%%%%%%%%%%%%%%%%%%%%%%%%%%%%%%%%%%%%%%%%%%%%%%

%%%%%%%%%%%%%%%% TRẮC NGHIỆM 4 PHƯƠNG ÁN %%%%%%%%%%%%%%%%

% nguồn: Lop_12_Dai_so_Chuong_3
\begin{ex}%[ma:ff530425]
	 %[2D3H1-3]
	Thời gian công tác của $50$  nhân viên một công ty được ghi lại ở bảng số liệu ghép nhóm sau (đơn vị năm).
	\begin{center}
		\begin{tabular}{|l|c|c|c|c|c|c|}
			\hline Nhóm  &{$[2;6)$}&{$[6;10)$}&{$[10;14)$}&{$[14;18)$}&{$[18;22)$}&{$[22;26)$}\\
			\hline Tần số & $10$ & $20$ &  $16$ & $3$ & $0$ & $1$\\
			 \hline
		\end{tabular}
	\end{center}
	Khoảng tứ phân vị của  mẫu số liệu ghép nhóm trên là
	\choice 
		{\True $5{,}375$} 
		 {$6{,}375$}
		{$8{,}875$} 
		 {$9{,}375$}
	\loigiai{Cỡ mẫu  $n=50$.\\
		Gọi $x_1; x_2; \ldots; x_{50}$ là mẫu số liệu gốc gồm thời gian công tác của $50$ nhân viên được xếp theo thứ  tự không giảm.\\
		Ta  có
		\begin{listEX}[3]
			\item $x_1, \ldots, x_{10}  \in[2;6)$
			 \item $x_{11}, \ldots, x_{30} \in[6;10)$
			\item $x_{31}, \ldots, x_{46}  \in[10;14)$
			\item  $x_{46}, \ldots, x_{49} \in[14;18)$
			\item  $x_{50} \in[22;26)$.
		\end{listEX}
		Tứ phân vị thứ nhất của mẫu số liệu gốc là $x_{13}  \in[6;10)$.\\
		Do đó, tứ phân vị  thứ nhất của mẫu số liệu ghép nhóm là
		\[Q_1=6+\dfrac{\dfrac{50}{4}-10}{20}  \cdot(10-6)=6{,}5.\]
		Tứ phân vị thứ ba  của mẫu số liệu gốc là $x_{38}$.\\ Mà $ x_{38} \in[10;14)$.\\ Do đó, tứ phân vị thứ ba của mẫu số liệu ghép nhóm là
		 \[Q_3=10+\dfrac{\dfrac{3\cdot50}{4}-30}{16}\cdot(14-10)=11{,}875.\]
		Vậy khoảng  tứ phân vị của mẫu số liệu ghép nhóm là
	 \[\Delta_Q=11,875-6{,}5=5{,}375.\]}
\end{ex}

% nguồn: Lop_12_Dai_so_Chuong_3
\begin{ex}%[ma:aa022f0d]
	 %[2D3H1-3]
	Bảng dưới đây thống kê độ ẩm không khí trung bình mỗi ngày đo được tại một trạm quan  trắc trong $60$ ngày đầu năm $2023$.
	\begin{center}
		\resizebox{0.8\linewidth}{!}{\begin{tikzpicture}
		\draw[opacity=.25,thin,step=.2,cyan](0,0) grid (7,2.5); 
		\draw[opacity=1,cyan](0,0) grid (7,2.5); 
		\draw[stealth-stealth](0,2.5)  node[left]{Tần số tương đối (\%)}|-(7,0) node[below,anchor=west]{Độ ẩm (\%)};
		\foreach\x/\dientich[count=\i from  1] in {0/60,1/64,1/68,1.5/72,1/76,0.5/80}{
			\draw[pattern=north  west lines,pattern color=cyan!65](\i-1,0) rectangle +(1,\x);
			\draw  (\i,0) node[below]{\dientich};
		} 
		\foreach  \y [count=\i from 1] in {10,20,30}{
			\draw (-.1,\i/2)--(.1,\i/2) 
			(0,\i/2)  node[left]{$\y$};}
		\foreach \z [count=\i from  1] in {20,20,30,20,10}{
			\draw (\i+0.5,\z/20)  node[above] {$\z$};}
	\end{tikzpicture}
		} 
	\end{center}
	 Tứ phân vị thứ nhất của mẫu số liệu ghép nhóm trên là
	 \choice
		{$64{,}25$} 
		{\True $65$} 
		{$66$} 
		 {$70$}
	\loigiai{Ta có bảng thống  kê theo nhóm như sau
		\begin{center}
			\begin{tabular}{|c|c|c|c|c|c|c|}
				 \hline
				Nhóm&  $[60;64)$ & $[64;68)$ & $[68;72)$ & $[72;76)$ & $[76;80)$ \\
				\hline 
				Tần số &$12$& $12$ & $18$  &$12$ & $6$ \\
				 \hline
			\end{tabular}
		\end{center}Cỡ mẫu $n=60$.\\
		Gọi $x_1; x_2; \ldots; x_{60}$ là mẫu số liệu gốc độ ẩm không khí trung bình mỗi ngày đo được tại một trạm quan trắc trong 60 ngày đầu năm 2023 được xếp theo  thứ tự không giảm.\\
		Ta  có
		\begin{listEX}[3]
			\item $x_1, \ldots, x_{12} \in[60;64)$ 
			\item $x_{13}, \ldots, x_{24}  \in[64;68)$
			\item $x_{25},  \ldots, x_{42} \in[68;72)$
			\item  $x_{43}, \ldots, x_{54} \in[72;76)$
			\item $x_{55}  \ldots, x_{60}\in[76;80)$.
		\end{listEX}
		Tứ phân vị  thứ nhất của mẫu số liệu gốc là $\dfrac{1}{2}(x_{15}+x_{16}) \in[64;68)$.\\
		Do đó, tứ phân vị thứ nhất của mẫu số liệu ghép  nhóm là
		\[Q_1=64+\dfrac{\dfrac{60}{4}-12}{12} \cdot(68-64)=65.\] 
	 }
\end{ex}

% nguồn: Lop_12_Dai_so_Chuong_3
\begin{ex}%[ma:d3e0d86c]
	 %[2D3H1-3]
	 Thời gian công tác của $50$ nhân viên một công ty được ghi lại ở bảng số liệu ghép nhóm sau (đơn vị năm).
	\begin{center}
		\begin{tabular}{|l|c|c|c|c|c|c|}
			\hline  Nhóm &{$[2;6)$}&{$[6;10)$}&{$[10;14)$}&{$[14;18)$}&{$[18;22)$}&{$[22;26)$}\\
			\hline Tần số & $10$ & $20$ & $16$ & $3$  & $0$ & $1$\\
			 \hline
		\end{tabular}
	\end{center}
	Tứ phân vị  thứ ba của mẫu số liệu ghép nhóm trên thuộc khoảng
	 \choice
		{$(5;7)$} 
		{$(16;17)$} 
		{\True  $(11;12)$}
		{$(15;16)$} 
	 \loigiai{Cỡ mẫu $n=50$.\\
		Gọi $x_1; x_2; \ldots; x_{50}$  là mẫu số liệu gốc gồm thời gian công tác của $50$ nhân viên được xếp theo thứ tự không giảm.\\
		Ta  có
		\begin{listEX}[3]
			\item $x_1, \ldots,  x_{10} \in[2;6)$
			\item $x_{11},  \ldots, x_{30} \in[6;10)$
			\item $x_{31}, \ldots, x_{46} \in[10;14)$ 
			\item $x_{46},  \ldots, x_{49} \in[14;18)$
			\item $x_{50} \in[22;26)$. 
		\end{listEX}
		Tứ phân vị thứ ba  của mẫu số liệu gốc là $\dfrac{1}{2}(x_{37}+x_{38})$.\\ Mà $x_{37}, x_{38} \in[10;14)$.\\ Do đó, tứ phân vị thứ ba của mẫu số liệu ghép nhóm là
		 \[Q_3=10+\dfrac{\dfrac{3\cdot50}{4}-30}{16}\cdot(14-10)=11,875.\]
	} 
\end{ex}

% nguồn: Lop_12_Dai_so_Chuong_3
\begin{ex}%[ma:7e806caf]
	%[2D3N1-2] 
	Tại một gara ô tô, người ta ghi  lại thời gian sử dụng của một số chiếc ô tô đến sữa chữa tại đây (đơn vị: năm). Kết quả được cho trong bảng sau đây.
	\begin{center}
		\begin{tabular}{|c|c|c|c|c|c|c|}
			\hline Nhóm & $[4; 6)$ & $[6; 8)$ & $[8; 10)$ & $[10;  12)$ & $[12; 14)$ & $[14; 16)$ \\
			\hline Tần số & $10$  & $14$ & $19$ & $23$ & $27$ & $32$ \\
			 \hline
		\end{tabular}
	\end{center}
	Khoảng biến thiên của mẫu số liệu ghép nhóm trên  là
	 \choice
		{$16$} 
		 {$14$}
		{\True $12$} 
		{$10$} 
	\loigiai{ 
		Khoảng biến thiên của mẫu số liệu  ghép nhóm trên là $R=16-4=12$.
	 }
\end{ex}

% nguồn: Lop_12_Dai_so_Chuong_3
\begin{ex}%[ma:c6cabb06]
	 %[2D3H2-2]
	Bảng sau ghi lại nhiệt độ không khí đo được lúc $12$ giờ trưa hằng ngày tại  một trạm quan sát trong $1$ năm.
	\begin{center}
		\begin{tabular}{|l|c|c|c|c|c|}
			\hline  Nhóm &{$[12;16)$}&{$[16;20)$}&{$[20;24)$}&{$[24;28)$}&{$[28;32)$}\\
			\hline Tần số & $85$ & $63$ & $95$ & $87$ & $35$ \\ 
			 \hline
		\end{tabular}
	\end{center}
	Độ lệch chuẩn của mẫu số  liệu ghép nhóm trên (làm tròn đến hàng phần nghìn) là
	\choice 
		{\True  $5{,}189$}
		 {$26{,}922$}
		 {$4{,}661$}
		{$21{,}727$} 
	\loigiai{Ta có bảng thống kê theo  giá trị đại diện như sau
		\begin{center}
			\begin{tabular}{|c|c|c|c|c|c|c|}
				\hline 
				Giá trị đại diện& $14$ & $18$ & $22$ &  $26$ & $30$ \\
				\hline 
				 Tần số &$85$& $63$ & $95$ &$87$ & $35$ \\
				\hline 
			\end{tabular}
		\end{center}
		Cỡ mẫu là $n=365$. Số trung bình của mẫu số liệu ghép nhóm trên là  \[\overline{x}=\dfrac{85\cdot 14+63\cdot 18+95\cdot 22+87\cdot 26+35\cdot 30}{365}\] \\
		Phương sai của  mẫu số liệu ghép nhóm là
		\[S^2=\dfrac{1}{365}\left(85\cdot 14^2+63\cdot 18^2+95\cdot 22^2+87\cdot 26^2+35\cdot  30^2\right)-\overline{x}^2\]
		Độ lệch chuẩn của  mẫu số liệu ghép nhóm là
	\[S =  \sqrt{S^2} \approx 5{,}189.\]}
\end{ex}

% nguồn: Lop_12_Dai_so_Chuong_3
\begin{ex}%[ma:f21956bc]
	%[2D3H2-2] 
	Điểm trung bình môn Toán của một số  học sinh khối $12$ được ghi lại ở bảng sau
	\begin{center}
		\begin{tabular}{|l|c|c|c|c|c|}
			\hline Nhóm  &{$[5;6)$}&{$[6;7)$}&{$[7;8)$}&{$[8;9)$}&{$[9;10)$}\\
			\hline Tần số  & $10$ & $14$ & $19$ & $32$ & $25$ \\
			\hline 
		\end{tabular}
	\end{center}
	Độ lệch chuẩn của mẫu số liệu ghép nhóm trên thuộc khoảng 
	 \choice
		{\True $(  1{,}25;1{,}28)$}
		{$(  1{,}28;1{,}3)$}
		 {$( 1{,}6;1{,}65)$}
		 {$( 1{,}65;1{,}7)$}
	\loigiai{Ta có bảng thống kê theo giá trị đại  diện như sau
		\begin{center}
			\begin{tabular}{|c|c|c|c|c|c|c|}
				 \hline
				 Điểm& $5{,}5$ & $6{,}5$ & $7{,}5$ & $8{,}5$ & $9{,}5$ \\
				 \hline
				Tần số &$10$& $14$  & $19$ &$32$ & $25$ \\
				\hline 
			\end{tabular}
		\end{center}
		Tổng số học sinh là $n=100$. Số trung bình  của mẫu số liệu ghép nhóm trên là \[\overline{x}=\dfrac{10\cdot 5{,}5+14\cdot 6{,}5+19\cdot 7{,}5+32\cdot 8{,}5+25\cdot 9{,}5}{100}=7{,}87.\] \\
		 Phương sai của mẫu số liệu ghép nhóm là
		\[S^2=\dfrac{1}{100}\left(10\cdot 5{,}5^2+14\cdot 6{,}5^2+19\cdot 7{,}5^2+32\cdot 8{,}5^2+25\cdot  9{,}5^2\right)-7{,}98^2 = 1{,}6296.\]
		Độ lệch chuẩn của mẫu số liệu  ghép nhóm là
	\[S =  \sqrt{1,6296} \approx 1{,}2766.\]}
\end{ex}


%%%%%%%%%%%%%%%% ĐÚNG/SAI %%%%%%%%%%%%%%%%

% nguồn: Lop_12_Dai_so_Chuong_3
\begin{ex}%[ma:35562135]
	%[2D3H2-2] 
	Kết quả kiểm tra môn Tiếng Anh (cùng đề) của học  sinh hai lớp 12A và 12B được cho lần lượt bởi mẫu số liệu ghép nhóm ở \textit{Bảng~6}, \textit{Bảng~7}.
	\choiceTF 
		{\True Số trung bình cộng của hai mẫu số liệu  trên bằng nhau}
		 {Độ lệch chuẩn của mẫu số liệu lớp 12A nhỏ hơn $2$}
		{\True  Phương sai của mẫu số liệu lớp 12B lớn hơn $3$}
		{\True Điểm thi của học sinh lớp 12B đồng đều hơn  lớp 12A}
		\begin{center}
			\begin{tabular}{|c|c|}
				 \hline
				Nhóm & Tần số \\ 
				 \hline
				 $[0;2)$& $3$\\
				\hline 
				 $[2;4)$&$5$\\
				\hline 
				 $[4;6)$&$5$\\
				\hline 
				$[6;8)$&$25$\\ 
				\hline 
				 $[8;10)$&$2$\\
				\hline 
				&$n=40$\\ 
				 \hline
		\end{tabular}
		 \hspace{0.5cm}
		\begin{tabular}{|c|c|}
			\hline 
			Nhóm &  Tần số \\
			\hline 
			 $[0;2)$& $1$\\
			 \hline
			 $[2;4)$&$4$\\
			\hline 
			 $[4;6)$&$15$\\
			 \hline
			 $[6;8)$&$16$\\
			\hline 
			 $[8;10)$&$4$\\
			\hline 
			&$n=40$\\ 
			\hline 
		\end{tabular} \\
		 \textit{Bảng 6} \hspace{2cm} \textit{Bảng 7}
	\end{center}
	 \loigiai{
		\begin{itemchoice}
			\itemch {\bf Đúng}.  \\
			 Ta có \textit{Bảng 6}
			\begin{center}
				\begin{tabular}{|c|c|c|c|c|c|}
					 \hline
					Nhóm &$[0;2)$&$[2;4)$&$[4;6)$&$[6;8)$&$[8;10)$\\ 
					 \hline
					Giá  trị đại diện & $1$&$3$&$5$&$7$&$9$\\
					 \hline
					Tần số&$3$&$5$&$5$&$25$&$2$\\ 
					\hline 
				\end{tabular}
			\end{center}
			Số trung bình của mẫu số liệu ghép nhóm lớp 12A là 
			\[ 
			\overline{x}=\dfrac{1}{40}\left(3\cdot 1+5\cdot 3+5\cdot 5+25\cdot  7+2\cdot 9\right)=5{,}9.
			 \]
			Ta có  \textit{Bảng 7}
			\begin{center}
				\begin{tabular}{|c|c|c|c|c|c|}
					\hline 
					Nhóm  &$[0;2)$&$[2;4)$&$[4;6)$&$[6;8)$&$[8;10)$\\
					\hline 
					Giá trị đại diện  & $1$&$3$&$5$&$7$&$9$\\
					\hline 
					Tần số&$1$&$4$&$15$&$16$&$4$\\ 
					 \hline
				\end{tabular}
			\end{center}
			 Số trung bình của mẫu số liệu ghép nhóm lớp 12B là
			 \[
			\overline{x}  = \dfrac{1}{40}\left(1\cdot 1+4\cdot 3+15\cdot 5+16\cdot 7+4\cdot 9\right) =5{,}9.
			\] 
			\itemch {\bf Sai}. \\ 
			Ta  có phương sai của mẫu số liệu ghép nhóm lớp 12A là
			 \allowdisplaybreaks
			\begin{eqnarray*}
				 s_A^2 &=& \dfrac{1}{40}\big[3(1-5{,}9)^2+5(3-5{,}9)^2+5(5-5{,}9)^2 \\
				&&+\, 25(7-5{,}9)^2+2(9-5{,}9)^2 
				 \big]=4{,}19.
			\end{eqnarray*}
			Độ lệch chuẩn mẫu số liệu ghép  nhóm lớp 12A là $s_A=\sqrt{s^2}=\sqrt{4{,}19}\approx 2{,}05$.\\
			\itemch {\bf Đúng}. \\ 
			Ta có  phương sai của mẫu số liệu ghép nhóm lớp $12B$ là
			\allowdisplaybreaks 
			\begin{eqnarray*}
				s_B^2 &=& \dfrac{1}{40}\big[1(1-5{,}9)^2+4(3-5{,}9)^2+15(5-5{,}9)^2  \\
				&&+\,16(7-5{,}9)^2+4(9-5{,}9)^2 
				 \big]=3{,}19.
			\end{eqnarray*}
			\itemch {\bf  Đúng}. \\
			Độ lệch chuẩn của lớp 12B là $\sqrt{3{,}19}\approx 1{,}79$. Do độ lệch chuẩn của lớp 12B nhỏ hơn độ lệch chuẩn của lớp 12A nên điểm thi lớp 12B đồng đều hơn. 
		\end{itemchoice}
	 }
\end{ex}


%%%%%%%%%%%%%%%% TRẢ LỜI NGẮN %%%%%%%%%%%%%%%%

% nguồn: Lop_12_Dai_so_Chuong_3
\begin{bt}%[ma:f0a14599]
	%[2D3V2-2] 
	\immini{ 
	Bảng bên biểu diễn mẫu số liệu ghép  nhóm về nhiệt độ không khí trung bình các tháng năm $2011$ tại Hà Nội (đơn vị: độ C) (\textit{Nguồn: Niên giám Thống kê $2021$, NXB Thống kê, $2022$}). Phương sai của mẫu số liệu đó bằng bao nhiêu (làm tròn kết quả đến hàng phần mười)?
	}{ 
	\begin{tabular}{|c|c|}
		 \hline
		 Nhóm & Tần số \\
		 \hline
		$[16{,}8;19{,}8)$ &$2$\\ 
		 \hline
		$[19{,}8;22{,}8)$  &$3$\\
		\hline 
		$[22{,}8;25{,}8)$ &$2$\\ 
		 \hline
		$[25{,}8;28{,}8)$ &$1$\\ 
		\hline 
		$[28{,}8;31{,}8)$ &$4$\\ 
		\hline 
		& $n=12$ \\ 
		 \hline
	\end{tabular}
	} 
	 \shortans{20,8}
	 \loigiai{
		 Ta có
		\begin{center}
			\resizebox{0.8\linewidth}{!}{\begin{tabular}{|c|c|c|c|c|c|}
			 \hline
			Nhóm & $[16{,}8;19{,}8)$& $[19{,}8;22{,}8)$  &$[22{,}8;25{,}8)$& $[25{,}8;28{,}8)$&$[28{,}8;31{,}8)$\\
			 \hline
			Giá trị đại diện  & $18{,}3$ &$21{,}3$&$24{,}3$&$27{,}3$&$30{,}3$\\
			 \hline
			Tần  số & $2$&$3$&$2$&$1$&$4$\\
			\hline 
		\end{tabular}}
	\end{center}
	Số trung bình của mẫu  số liệu ghép nhóm là
	 \allowdisplaybreaks
	\begin{eqnarray*}
		\overline{x}=\dfrac{1}{12}\left(2\cdot 18{,}3 +3\cdot  21{,}3+2\cdot 24{,}3+1\cdot 27{,}3+4\cdot 30{,}3\right) = 24{,}8.
	\end{eqnarray*}
	Phương sai của mẫu số liệu  ghép nhóm
	\allowdisplaybreaks 
	\begin{eqnarray*}
		 s^2 &=& \dfrac{1}{12}\big[2(18{,}3-24{,}8)^2+3(21{,}3-24{,}8)^2+2(24{,}3-24{,}8)^2 \\
		 &&+\,1(27{,}3-24{,}8)^2
		+4(30{,}3-24{,}8)^2 
		\big]=\dfrac{83}{4}=20{,}75. 
	\end{eqnarray*}
 }
\end{bt}

% nguồn: Lop_12_Dai_so_Chuong_3
\begin{bt}%[ma:e45ee6f7]
	 %[2D3H1-1]
	Biết hàm số $F(x)=a \sin x+b \cos x$ là một nguyên hàm của  hàm số $f(x)=2 \sin x-\cos x$. Khi đó giá trị của $a+2 b$ bằng
	 \shortans{-5}
	 \loigiai{
		 Ta có
		\[ 
		\displaystyle\int f(x)\mathrm{\,d}x 
		=\displaystyle\int(2  \sin x-\cos x)\mathrm{\,d}x
		=-2  \cos x-\sin x+C.
		\] 
		Suy ra $F(x)=-\sin x-2 \cos  x$ hay $a=-1$ và $b=-2$.\\
		Do đó $a+2  b=-5$.
	 }
\end{bt}

% nguồn: Lop_12_Dai_so_Chuong_3
\begin{bt}%[ma:a5adf192]
	 %[2D3H2-2]
	Bảng sau  đây biểu diễn mẫu số liệu ghép nhóm về cân nặng của một số quả dưa bở thu hoạch được ở một khu vườn (đơn vị: gam).
	\begin{center}
		\begin{tabular}{|l|c|c|c|c|c|}
			\hline  Nhóm &{$[600; 650)$}&{$[650; 700)$}&{$[700; 750)$}&{$[750; 800)$}&{$[800; 850)$}\\
			\hline Tần số  & $14$ & $40$ & $13$ & $10$ & $3$ \\
			 \hline
		\end{tabular}
	\end{center}
	Tìm phương sai của mẫu số liệu ghép nhóm trên (kết quả làm  tròn đến hàng đơn vị).
	\shortans{2631} 
	\loigiai{Trung bình của mẫu số  liệu \[\overline{x}=\dfrac{14\cdot 625+40\cdot 675+13\cdot 725+10\cdot 775+3\cdot 825}{14+40+13+10+3}=\dfrac{1\,385}{2}.\]
		Phương sai là 
		\begin{eqnarray*}
			&s^2&=\dfrac{14\cdot  (625-\overline{x})^2+40\cdot (675-\overline{x})^2+13\cdot (725-\overline{x})^2+10\cdot (775-\overline{x})^2+3\cdot (825-\overline{x})^2}{80}\\
			 &&\approx 2\,631.
		\end{eqnarray*}
	 }
\end{bt}

% nguồn: Lop_12_Dai_so_Chuong_3
\begin{bt}%[ma:d5c29042]
	%[2D3V2-3] 
	Bảng tần số tương đối ghép nhóm sau đây ghi lại huyết áp tâm thu  của 80 người cao tuổi ở một khu vực (đơn vị mmHg)
	\begin{center}
		\begin{tabular}{|l|c|c|c|c|c|}
			\hline Nhóm  &{$[120;125)$}&{$[125;130)$}&{$[130;135)$}&{$[135;140)$}&{$[140;145)$}\\
			 \hline Tần số & $5\% $ & $10\%$ & $15\%$ & $25\%$ & $45\%$ \\
			 \hline
		\end{tabular}
	\end{center}
	Tìm mốt của mẫu số liệu ghép nhóm trên (kết quả làm tròn đến hàng đơn  vị).
	 \shortans {$142$}
	\loigiai{Bảng tần số ghép nhóm\begin{center}
		\begin{tabular}{|l|c|c|c|c|c|}
			\hline Nhóm  &{$[120;125)$}&{$[125;130)$}&{$[130;135)$}&{$[135;140)$}&{$[140;145)$}\\
			\hline Tần số & $4 $ &  $8$ & $12$ & $20$ & $36$ \\
			 \hline
		\end{tabular}
	\end{center}
	Nhóm chứa mốt là  $[140;145)$\\
	Mốt của mẫu số liệu ghép  nhóm là
 \[M_0=140+\dfrac{36-20}{(36-20)+(36-0)}\cdot(145-140)\approx141{,}54\]}
\end{bt}

% nguồn: Lop_12_Dai_so_Chuong_3
\begin{bt}%[ma:4d49c486]
	%[2D3H2-2] 
	Bảng dưới đây thống  kê cự li ném tạ của một vận động viên.\\
	\centerline{\begin{tabular}{|c|c|c|c|c|c|}
	\hline 
	Cự  li (m) &[19; 19,5) &[19,5; 20) &[20; 20,5) &[20,5; 21) &[21; 21,5)\\
	\hline 
	Tần số & 13 &  45 & 24 & 12 & 6\\
	\hline 
\end{tabular}}
Tính phương sai của mẫu số liệu ghép nhóm trên. (Kết quả làm tròn  đến hàng phần trăm)
\shortans{0,28} 
\loigiai{ 
	 Có bảng giá trị đại diện:\\
	\centerline{\begin{tabular}{|c|c|c|c|c|c|}
	\hline 
	GTĐD & 19,25 & 19,75 & 20,25  & 20,75 & 21,25\\
	\hline 
	Tần số & 13 & 45 & 24 & 12 &  6\\
	 \hline
	\end{tabular}}
	Giá trị  trung bình:\\
	$\overline{x}=\dfrac{13.19,25+45.19,75+24.20,25+12.20,75+6.21,25}{100}=20,015$\\ 
	Phương sai của mẫu số liệu:\\ 
 $s^2=\dfrac{13.19,25^2+45.19,75^2+24.20,25^2+12.20,75^2+6.21,25^2}{100}-20,015^2\approx 0,28$.}
\end{bt}
